\section{Intern analyse}

Ifølge ressursbasert teori springer konkurransefortrinn ut av bedriftens
unike interne ressurser og kapabiliteter snarere enn av markedsstrukturer
alene \parencite{barney1991}. Med dette som teoretisk utgangspunkt vil
internanalysen vurdere hvilke styrker og svakheter Zipline har i møte med en
potensiell etablering i det norske dagligvare- og restaurantmarkedet.
Sentralt i analysen står hvilket verdiforslag selskapet faktisk tilbyr,
hvilke ressurser og kapabiliteter som understøtter dette, hvor krevende de
er å imitere for konkurrenter, og hvilke begrensninger som kan svekke
tjenestens konkurransekraft.

\subsection{Styrker}

En av Ziplines mest fundamentale styrker er selskapets teknologiske og
operasjonelle kompetanse. Den avanserte droneteknologien og den særpregede
leveringsmetoden skiller Zipline tydelig fra tradisjonelle
leveringsalternativer. I tillegg har selskapet opparbeidet betydelig
operasjonell erfaring gjennom over 2,1 millioner autonome leveranser siden
2016 \parencite{ziplineabout}. Størstedelen av dette volumet stammer
imidlertid fra medisinske leveranser i afrikanske markeder, mens den
forbrukerrettede dagligvare- og restauranttjenesten foreløpig kun er
etablert i Dallas-Fort Worth og Pea Ridge \parencite{rogers2024}. Erfaringen
gir dermed Zipline reell modenhet innen autonom drift, flåtestyring og
regulatorisk arbeid, men direkte overførbarhet til det norske dagligvare- og
restaurantsegmentet er mer begrenset enn det samlede volumtallet alene
antyder. Et sentralt forbehold er i tillegg at både teknologien og
erfaringsgrunnlaget primært er opparbeidet i operative omgivelser som
skiller seg vesentlig fra det norske klimaet, noe som behandles nærmere
under svakheter.

En ytterligere styrke ved Zipline er at selve leveringsmodellen er bygget
rundt autonom drift. I motsetning til tradisjonelle hjemleveringsaktører
som Foodora og Wolt, der hver enkelt levering avhenger av en menneskelig
sjåfør, erstatter Zipline store deler av denne arbeidskostnaden med autonom
drift. Ziplines egne estimater antyder at enhetskostnaden ved Platform
2-leveranser kan falle til 2 til 4 dollar per leveranse ved modenhet
\parencite{logisticsnavigators2026}. En slik kostnadsstruktur er vanskelig å
imitere for eksisterende aktører uten å bygge om hele verdikjeden, ettersom
forretningsmodellene deres er innrettet rundt manuell arbeidskraft. I et
norsk marked, der lønnskostnadene ligger blant de høyeste i Europa
\parencite{ssb2024}, vil en slik strukturell forskjell mellom
arbeidskraftbaserte og autonome modeller ha større praktisk betydning enn i
lavkostmarkeder.

Leveringstiden er også en sentral intern ressurs ved plattformen. Som nevnt
i produktbeskrivelsen har Zipline en gjennomsnittlig leveringstid på 15
minutter. Til sammenlikning oppgir Foodora selv en gjennomsnittlig
leveringstid på 30 minutter \parencite{foodora_hahalden}, og Ziplines
leveringstempo ligger dermed klart under det etablerte aktører i det norske
dagligvare- og restaurantmarkedet opererer med i dag. Tempoet er en direkte konsekvens av den
dronebaserte arkitekturen, som unngår veinettet og dermed både
trafikkforsinkelser og manuelle stopp. Det kan derfor ikke enkelt
replikeres av aktører hvis siste ledd er basert på bil, sykkel eller
gange, uten at hele distribusjonsmodellen legges om.

Ziplines løsning har også en distinkt bærekraftsprofil. Ifølge selskapet
reduseres utslipp med 97\,\% sammenliknet med tradisjonelle
leveringsmetoder \parencite{rogers2024}, og i en norsk kontekst forsterkes
egenskapen av at den nasjonale strømmiksen er tilnærmet helelektrisk
fornybar. Profilen er likevel ikke enestående i markedet, ettersom både
elektrifiserte budbiler og sykkelbud allerede viser lave utslipp i egne
segmenter. Forspranget blir dessuten mindre når konkurrentene elektrifiserer flåten.
Zipline bruker likevel bærekraft aktivt i både produktdesign og
markedsføring.

\subsection{Svakheter}

En sentral svakhet ved Zipline er at tjenesten har et relativt begrenset
bruksområde. Som nevnt i produktbeskrivelsen har Platform 2, som er dronen
rettet mot hjemlevering, en kapasitet på 3,6 kg og en leveringsradius på 16
km. Disse begrensningene gjør at tjenesten først og fremst passer til
mindre og relativt enkle leveranser, og at mange kundebehov faller utenfor
det plattformen kan dekke. Ressursen er dermed verdifull i et smalere
utsnitt av hjemleveringsmarkedet enn tradisjonelle leveringsalternativer.

I tillegg har Zipline begrenset dokumentert erfaring med operativ drift
under norske vinterforhold. Selskapet har primært operert i
Afrika og USA, og det er usikkert i hvilken grad den eksisterende
teknologiske plattformen er testet og tilpasset forhold med snø, kulde og
lange perioder med mørketid. Til sammenligning har den norske aktøren
Aviant demonstrert operasjonell drift ned til minus 26 grader i Midt-Norge
\parencite{aviant2024}. Svakheten gjelder dermed både teknologiens
robusthet og selskapets driftskompetanse i arktiske omgivelser.

En annen svakhet er at Ziplines modell er operasjonelt og kommersielt
kompleks. Løsningen forutsetter ikke bare avansert teknologi, men også tett
koordinering mellom drift, leveringspunkter, partnerintegrasjoner og
kundegrensesnitt. Å bygge opp slike koordineringskapabiliteter på nytt for
hvert marked er krevende, og Zipline har foreløpig ikke dokumentert at
modellen fungerer utenfor sine etablerte operasjonsområder. Svakheten
handler dermed ikke om selve teknologien, men om organisasjonens evne til å
operasjonalisere den i en ny kontekst.


% dette avsnittet burde integreres inn i SWOT
En ytterligere svakhet ved Zipline er utfordringen knyttet til å oppnå en
gunstig balanse mellom customer acquisition cost (CAC) og customer lifetime
value (LTV). CAC kan bli høy fordi Zipline ikke bare introduserer en ny
tjeneste, men også en ny leveringsform som kan kreve betydelige ressurser til
markedsføring, opplæring og tillitsbygging. At tjenesten bestilles gjennom
Ziplines egen app, kan også skape ekstra friksjon i anskaffelsesfasen, siden
kundene må ta i bruk en ny plattform fremfor å benytte løsninger de allerede
kjenner. Samtidig kan LTV bli begrenset fordi tjenesten virker mest relevant i
bestemte og situasjonsavhengige kjøp, særlig mindre og tidskritiske
leveranser. Begrensninger i leveringsvekt, radius og varetyper innebærer dessuten
at tjenesten i praksis bare dekker en del av kundens samlede behov for
hjemlevering. Dette kan føre til at mange kunder bare bruker tjenesten sporadisk, 
noe som gjør det vanskeligere å bygge høy bruksfrekvens og sterk kundelojalitet over tid.
